\documentclass[11pt,a4paper]{article}

% ==================== TEMEL PAKETLER (xelatex) ====================
\usepackage{fontspec}
\usepackage{polyglossia}
\setmainlanguage{turkish}
\usepackage[a4paper,margin=2cm,top=2.4cm,bottom=2.4cm]{geometry}
\usepackage{xcolor}
\usepackage{graphicx}
\usepackage{array}
\usepackage{booktabs}
\usepackage{longtable}
\usepackage{tabularx}
\usepackage{colortbl}
\usepackage{multirow}
\usepackage{enumitem}
\usepackage[most]{tcolorbox}
\usepackage{listings}
\usepackage{fancyhdr}
\usepackage{titlesec}
\usepackage{titletoc}
\usepackage{hyperref}
\usepackage{fontawesome5}
\setlength{\headheight}{13.6pt}

% ==================== FONTLAR ====================
\setmainfont{Latin Modern Roman}
\newfontfamily\codefont{Consolas}[Scale=0.88]

% ==================== RENK PALETİ ====================
\definecolor{anaMavi}{HTML}{1B3A5B}
\definecolor{vurguMavi}{HTML}{2E6DA4}
\definecolor{acikMavi}{HTML}{EAF2F9}
\definecolor{turuncu}{HTML}{D9701E}
\definecolor{yesil}{HTML}{2E7D32}
\definecolor{acikYesil}{HTML}{E8F5E9}
\definecolor{kirmizi}{HTML}{C0392B}
\definecolor{mor}{HTML}{6C3483}
\definecolor{acikMor}{HTML}{F3E9F7}
\definecolor{acikGri}{HTML}{F4F5F7}
\definecolor{ortaGri}{HTML}{6B7280}
\definecolor{koyuGri}{HTML}{2D2D2D}
\definecolor{kodArka}{HTML}{FBFBFA}
\definecolor{satirRenk}{HTML}{F0F4F8}
\definecolor{qtYesil}{HTML}{41CD52}

% ==================== HYPERREF ====================
\hypersetup{
  colorlinks=true,
  linkcolor=vurguMavi,
  urlcolor=vurguMavi,
  pdftitle={QML ve C++ Entegrasyonu Kapsamlı Rehber},
  pdfauthor={QML C++ Rehberi},
  bookmarksnumbered=true
}

% ==================== BAŞLIK STİLLERİ ====================
\titleformat{\section}
  {\color{anaMavi}\Large\bfseries}
  {\thesection.}{0.6em}{}
  [\vspace{2pt}{\color{turuncu}\titlerule[1.6pt]}]
\titleformat{\subsection}
  {\color{vurguMavi}\large\bfseries}
  {\thesubsection}{0.6em}{}
\titleformat{\subsubsection}
  {\color{koyuGri}\normalsize\bfseries}
  {\thesubsubsection}{0.6em}{}
\titlespacing*{\section}{0pt}{16pt}{8pt}

% ==================== BAŞLIK / ALTBİLGİ ====================
\pagestyle{fancy}
\fancyhf{}
\renewcommand{\headrulewidth}{0.4pt}
\renewcommand{\footrulewidth}{0.4pt}
\fancyhead[L]{\small\color{ortaGri}QML \& C++ Entegrasyon Rehberi}
\fancyhead[R]{\small\color{ortaGri}\leftmark}
\fancyfoot[C]{\small\color{ortaGri}\thepage}

% ==================== KOD LİSTELEME: C++ ====================
\lstdefinestyle{cppstyle}{
  language=C++,
  basicstyle=\codefont\small\color{koyuGri},
  keywordstyle=\color{vurguMavi}\bfseries,
  commentstyle=\color{yesil}\itshape,
  stringstyle=\color{turuncu},
  numberstyle=\tiny\color{ortaGri},
  numbers=left,
  numbersep=8pt,
  backgroundcolor=\color{kodArka},
  frame=single,
  framesep=6pt,
  rulecolor=\color{acikGri!80!ortaGri},
  breaklines=true,
  showstringspaces=false,
  tabsize=4,
  columns=fullflexible,
  keepspaces=true,
  morekeywords={Q_OBJECT,Q_PROPERTY,Q_INVOKABLE,Q_ENUM,Q_FLAG,Q_GADGET,
    QML_ELEMENT,QML_SINGLETON,QML_UNCREATABLE,QML_NAMED_ELEMENT,QML_ANONYMOUS,
    QML_FOREIGN,QML_ATTACHED,QML_VALUE_TYPE,signals,slots,Q_SIGNALS,Q_SLOTS,
    emit,Q_EMIT,override,nullptr,constexpr,explicit,final,qmlRegisterType,
    qmlRegisterSingletonType,setContextProperty,connect,MEMBER,NOTIFY,READ,
    WRITE,CONSTANT,REQUIRED,FINAL,BINDABLE,qRegisterMetaType}
}

% ==================== KOD LİSTELEME: QML ====================
\lstdefinelanguage{QML}{
  morekeywords={import,property,signal,function,id,anchors,width,height,color,
    Rectangle,Item,Text,MouseArea,Component,Window,ApplicationWindow,Column,Row,
    ListView,ListModel,Repeater,Button,TextField,Connections,readonly,alias,
    required,var,int,real,string,bool,true,false,null,on,Behavior,states,State,
    Timer,Loader,pragma,QtQuick,onClicked,onCompleted,model,delegate},
  sensitive=true,
  morecomment=[l]{//},
  morecomment=[s]{/*}{*/},
  morestring=[b]",
  morestring=[b]'
}
\lstdefinestyle{qmlstyle}{
  language=QML,
  basicstyle=\codefont\small\color{koyuGri},
  keywordstyle=\color{mor}\bfseries,
  commentstyle=\color{yesil}\itshape,
  stringstyle=\color{turuncu},
  numberstyle=\tiny\color{ortaGri},
  numbers=left,
  numbersep=8pt,
  backgroundcolor=\color{qtYesil!6},
  frame=single,
  framesep=6pt,
  rulecolor=\color{qtYesil!50!ortaGri},
  breaklines=true,
  showstringspaces=false,
  tabsize=4,
  columns=fullflexible,
  keepspaces=true
}
\lstset{style=cppstyle}

% ==================== ÖZEL KUTULAR ====================
\newtcolorbox{bilgikutu}[1][Bilgi]{
  colback=acikMavi, colframe=vurguMavi, boxrule=0.8pt, arc=3pt,
  left=8pt,right=8pt,top=6pt,bottom=6pt,
  fonttitle=\bfseries\color{white}, coltitle=white,
  title={\faInfoCircle~~#1}
}
\newtcolorbox{ipucukutu}[1][İpucu]{
  colback=acikYesil, colframe=yesil, boxrule=0.8pt, arc=3pt,
  left=8pt,right=8pt,top=6pt,bottom=6pt,
  fonttitle=\bfseries\color{white}, coltitle=white,
  title={\faLightbulb~~#1}
}
\newtcolorbox{uyarikutu}[1][Dikkat]{
  colback=kirmizi!8, colframe=kirmizi, boxrule=0.8pt, arc=3pt,
  left=8pt,right=8pt,top=6pt,bottom=6pt,
  fonttitle=\bfseries\color{white}, coltitle=white,
  title={\faExclamationTriangle~~#1}
}
\newtcolorbox{ornekkutu}[1][Örnek]{
  colback=acikGri, colframe=ortaGri!60, boxrule=0.7pt, arc=2pt,
  left=8pt,right=8pt,top=6pt,bottom=6pt,
  fonttitle=\bfseries\color{koyuGri}, coltitle=koyuGri,
  title={\faCode~~#1}
}
\newtcolorbox{notkutu}[1][Not]{
  colback=acikMor, colframe=mor, boxrule=0.8pt, arc=3pt,
  left=8pt,right=8pt,top=6pt,bottom=6pt,
  fonttitle=\bfseries\color{white}, coltitle=white,
  title={\faStickyNote~~#1}
}

% Yardımcı komutlar
\newcommand{\kls}[1]{\texttt{\color{anaMavi}\textbf{#1}}}
\newcommand{\hx}[1]{\texttt{\color{turuncu}#1}}
\newcommand{\kod}[1]{\texttt{\color{koyuGri}#1}}

\renewcommand{\arraystretch}{1.35}

% ==================== BAŞLANGIÇ ====================
\begin{document}

% ---------- KAPAK ----------
\begin{titlepage}
\centering
\vspace*{1.6cm}
{\color{qtYesil}\rule{\textwidth}{2pt}}\\[6pt]
{\color{anaMavi}\Huge\bfseries QML \& C++ Entegrasyonu}\\[10pt]
{\color{vurguMavi}\LARGE Kapsamlı Uygulama Rehberi}\\[8pt]
{\color{ortaGri}\large Qt~6 \textbullet\ Modern CMake \textbullet\ Bol Çalışan Örnek}\\[6pt]
{\color{qtYesil}\rule{\textwidth}{2pt}}\\[1.6cm]

\begin{tcolorbox}[width=0.88\textwidth,colback=acikMavi,colframe=vurguMavi,arc=4pt,boxrule=1pt]
\centering\large
Bu belge, bir Qt uygulamasında \textbf{C++ arka ucu} ile \textbf{QML arayüzü}
arasındaki köprüyü baştan sona anlatır: tip kaydı, özellikler (\kod{Q\_PROPERTY}),
çağrılabilir metotlar, sinyal--slot bağları, enum'lar, tekil nesneler (singleton),
liste modelleri ve iki yönlü iletişim. Her başlık \textbf{çok sayıda çalışan kod
örneğiyle} desteklenmiştir.
\end{tcolorbox}

\vspace{0.9cm}
\begin{tcolorbox}[width=0.88\textwidth,colback=acikGri,colframe=ortaGri!50,arc=3pt,boxrule=0.6pt]
\small
\textbf{Kapsam:} Meta-Object sistemi \textbullet\ CMake \& \kod{qt\_add\_qml\_module}
\textbullet\ \kod{QML\_ELEMENT} ile tip kaydı \textbullet\ \kod{Q\_PROPERTY} \&
NOTIFY \textbullet\ \kod{Q\_INVOKABLE} \& slotlar \textbullet\ sinyaller ve
\kod{Connections} \textbullet\ \kod{Q\_ENUM} \textbullet\ context property
\textbullet\ singleton'lar \textbullet\ \kod{QAbstractListModel} \textbullet\
C++'tan QML'e erişim \textbullet\ değer tipleri \textbullet\ yaygın hatalar
\end{tcolorbox}

\vfill
{\color{ortaGri}\large Hazırlanma Tarihi: 9 Temmuz 2026}\\[4pt]
{\color{ortaGri} Qt 6.x \textbullet\ C++17 \textbullet\ QtQuick}
\vspace*{0.6cm}
\end{titlepage}

% ---------- İÇİNDEKİLER ----------
\tableofcontents
\thispagestyle{empty}
\clearpage

% =====================================================================
\section{Giriş: QML ile C++ Neden Birlikte Kullanılır?}
% =====================================================================

Qt Quick uygulamalarında görev dağılımı nettir: \textbf{QML} akıcı, animasyonlu
ve deklaratif kullanıcı arayüzünü tanımlar; \textbf{C++} ise iş mantığını, veri
işlemeyi, dosya/ağ erişimini ve performans kritik hesaplamaları üstlenir. İkisini
birbirine bağlayan tutkal, Qt'nin \emph{Meta-Object} sistemidir.

\begin{bilgikutu}
\textbf{Neden ayrım?} QML yorumlanan (JavaScript motoru üzerinde çalışan) bir
dildir; hızlı prototipleme ve arayüz için idealdir ama ağır hesaplama için
uygun değildir. C++ ise derlenir ve çok hızlıdır. Doğru mimari: \emph{ince}
bir QML arayüz katmanı + \emph{kalın} bir C++ mantık katmanı. Bu belge, bu iki
katmanın nasıl konuşturulacağını gösterir.
\end{bilgikutu}

\subsection{Meta-Object Sistemi: Köprünün Temeli}

C++ normalde derleme sonrası tip bilgisini kaybeder (\emph{static typing}). QML'in
çalışma zamanında bir C++ nesnesinin özelliklerini okuması, metotlarını çağırması
ve sinyallerine bağlanması için \textbf{çalışma zamanı tip bilgisine (introspection)}
ihtiyaç vardır. Qt bunu \kls{QObject} taban sınıfı ve \kls{moc} (Meta-Object
Compiler) ile sağlar.

\begin{center}
\begin{tabularx}{\textwidth}{@{}>{\bfseries\color{anaMavi}}l >{\raggedright\arraybackslash}X@{}}
\toprule
\rowcolor{acikMavi} \color{anaMavi}\textbf{Bileşen} & \color{anaMavi}\textbf{Görevi} \\
\midrule
QObject & Sinyal/slot, özellik ve nesne ağacının temel taban sınıfı. QML'e açılacak her tip bundan türer. \\
\rowcolor{acikGri} Q\_OBJECT makrosu & Sınıfa meta-object desteği ekler. moc bu makroyu görünce ek kod üretir. \\
moc & Başlık dosyalarını tarar, \kod{Q\_OBJECT} içeren sınıflar için \kod{moc\_*.cpp} üretir. CMake bunu otomatik çalıştırır. \\
\rowcolor{acikGri} QML Engine & QML dosyalarını okur, JavaScript'i çalıştırır ve meta-object üzerinden C++ ile konuşur. \\
\bottomrule
\end{tabularx}
\end{center}

\begin{uyarikutu}
QML'e açtığınız her C++ sınıfı üç kurala uymalıdır: \textbf{(1)} \kls{QObject}'ten
(veya türevinden) türemeli, \textbf{(2)} sınıfın en başında \kod{Q\_OBJECT}
makrosu bulunmalı, \textbf{(3)} varsayılan (argümansız) kurucusu olmalı veya
tekil/oluşturulamaz olarak işaretlenmeli. Bu üçü olmadan QML tarafı sınıfı
göremez.
\end{uyarikutu}

\subsection{İki Yönlü İletişimin Dört Yolu}

C++ ile QML arasındaki bilgi akışı dört mekanizmayla olur. Bu belgenin
tamamı bu dört yolu derinleştirir:

\begin{center}
\begin{tabularx}{\textwidth}{@{}>{\bfseries\color{vurguMavi}}l >{\raggedright\arraybackslash}X c@{}}
\toprule
\rowcolor{acikMavi}\color{anaMavi}\textbf{Mekanizma} & \color{anaMavi}\textbf{Ne işe yarar} & \color{anaMavi}\textbf{Yön} \\
\midrule
Q\_PROPERTY & Veri açar; QML okur/yazar, değişince otomatik güncellenir & C++ $\leftrightarrow$ QML \\
\rowcolor{acikGri} Q\_INVOKABLE / slot & C++ metodunu QML'den çağırma & QML $\rightarrow$ C++ \\
Sinyal (signal) & C++'ta olay olunca QML'i haberdar etme & C++ $\rightarrow$ QML \\
\rowcolor{acikGri} Metot/özellik erişimi & C++'tan QML nesnesine erişip metot çağırma & C++ $\rightarrow$ QML \\
\bottomrule
\end{tabularx}
\end{center}

% =====================================================================
\section{Proje Kurulumu: CMake ve QML Modülü}
% =====================================================================

Qt 6'da önerilen yol, C++ tiplerini bir \textbf{QML modülü} içinde toplamaktır.
\kod{qt\_add\_qml\_module} komutu, tip kaydını otomatikleştirir ve QML derleyici
(qmlcachegen, qmltc) optimizasyonlarını mümkün kılar.

\subsection{Klasör Yapısı}

\begin{lstlisting}[language=bash,basicstyle=\codefont\small,keywordstyle=\color{koyuGri}]
MyApp/
├── CMakeLists.txt
├── main.cpp
├── backend.h          # C++ sınıfı (başlık)
├── backend.cpp        # C++ sınıfı (uygulama)
└── Main.qml           # Arayüz
\end{lstlisting}

\subsection{CMakeLists.txt — Tam Örnek}

\begin{lstlisting}[language=bash,morekeywords={cmake_minimum_required,project,find_package,qt_standard_project_setup,qt_add_executable,qt_add_qml_module,target_link_libraries,VERSION,COMPONENTS,REQUIRED,PRIVATE},keywordstyle=\color{vurguMavi}\bfseries]
cmake_minimum_required(VERSION 3.16)
project(MyApp VERSION 1.0 LANGUAGES CXX)

# Qt 6 Quick modulunu bul
find_package(Qt6 6.5 REQUIRED COMPONENTS Quick)

# AUTOMOC/AUTORCC ve C++ standardini ayarlar
qt_standard_project_setup(REQUIRES 6.5)

# Yurutulebilir hedef
qt_add_executable(appMyApp main.cpp)

# QML modulu: URI = "MyApp", surum = 1.0
qt_add_qml_module(appMyApp
    URI MyApp
    VERSION 1.0
    QML_FILES
        Main.qml
    SOURCES
        backend.h
        backend.cpp
)

target_link_libraries(appMyApp PRIVATE Qt6::Quick)
\end{lstlisting}

\begin{notkutu}[URI nedir?]
\kod{URI MyApp} satırı, QML tarafında \kod{import MyApp} yazarak bu modüldeki
tüm C++ tiplerine erişebileceğiniz anlamına gelir. Modül adı ile \kod{import}
ifadesi \textbf{birebir aynı} olmalıdır.
\end{notkutu}

\subsection{main.cpp — Motoru Başlatma}

\begin{lstlisting}
#include <QGuiApplication>
#include <QQmlApplicationEngine>

int main(int argc, char *argv[])
{
    QGuiApplication app(argc, argv);

    QQmlApplicationEngine engine;

    // Modul icindeki ana QML dosyasini yukle:
    // "MyApp" modulu + "Main" tipi
    engine.loadFromModule("MyApp", "Main");

    if (engine.rootObjects().isEmpty())
        return -1;   // QML yuklenemedi

    return app.exec();
}
\end{lstlisting}

\begin{ipucukutu}
\kod{loadFromModule("MyApp", "Main")} — Qt 6.5+ ile gelen modern yükleme
biçimidir. Eski \kod{engine.load(QUrl("qrc:/..."))} yönteminin yerini alır;
dosya yollarıyla uğraşmanıza gerek kalmaz.
\end{ipucukutu}

% =====================================================================
\section{C++ Tiplerini QML'e Kaydetme}
% =====================================================================

Bir C++ sınıfını QML'de kullanabilmek için önce onu \emph{kaydetmek} gerekir.
Qt 6'da bu, başlık dosyasına konan basit makrolarla yapılır; \kls{moc} ve
\kod{qt\_add\_qml\_module} gerisini halleder.

\subsection{QML\_ELEMENT — En Yaygın Yöntem}

Sınıfa \kod{QML\_ELEMENT} makrosunu ekleyin; sınıf, C++'taki adıyla QML'de
kullanılabilir hâle gelir.

\begin{ornekkutu}[backend.h — Kaydedilebilir bir sınıf]
\end{ornekkutu}
\begin{lstlisting}
#ifndef BACKEND_H
#define BACKEND_H

#include <QObject>
#include <QQmlEngine>   // QML_ELEMENT icin gerekli

class Backend : public QObject
{
    Q_OBJECT
    QML_ELEMENT          // <-- Bu sinifi QML'e "Backend" olarak kaydeder

public:
    // QML'in nesne olusturabilmesi icin argumansiz kurucu sart
    explicit Backend(QObject *parent = nullptr);

    Q_INVOKABLE QString greet(const QString &name) const;
};

#endif
\end{lstlisting}

\begin{lstlisting}[caption={}]
// backend.cpp
#include "backend.h"

Backend::Backend(QObject *parent) : QObject(parent) {}

QString Backend::greet(const QString &name) const
{
    return QStringLiteral("Merhaba, %1!").arg(name);
}
\end{lstlisting}

Artık QML tarafında hiçbir ek işlem olmadan kullanılabilir:

\begin{lstlisting}[style=qmlstyle]
import QtQuick
import MyApp          // CMake'teki URI ile ayni

Window {
    width: 300; height: 200; visible: true

    Backend { id: backend }   // C++ sinifindan nesne olusturuldu!

    Text {
        anchors.centerIn: parent
        text: backend.greet("Dunya")   // -> "Merhaba, Dunya!"
    }
}
\end{lstlisting}

\begin{bilgikutu}
\kod{QML\_ELEMENT}, sınıfı \textbf{C++ ismiyle} kaydeder. Farklı bir isim
istiyorsanız \kod{QML\_NAMED\_ELEMENT(BaskaIsim)} kullanın. Kaydın çalışması için
sınıfın bir QML modülünün parçası olması (yani \kod{SOURCES} altında listelenmesi)
yeterlidir; ayrı bir \kod{qmlRegisterType} çağrısına gerek yoktur.
\end{bilgikutu}

\subsection{Kayıt Makroları Karşılaştırması}

\begin{center}
\begin{tabularx}{\textwidth}{@{}>{\ttfamily\color{anaMavi}\small}l >{\raggedright\arraybackslash}X@{}}
\toprule
\rowcolor{acikMavi}\color{anaMavi}\textbf{Makro} & \color{anaMavi}\textbf{Ne zaman kullanılır} \\
\midrule
QML\_ELEMENT & Sınıfı C++ adıyla açar. Standart, en sık kullanılan. \\
\rowcolor{acikGri} QML\_NAMED\_ELEMENT(Ad) & QML'de farklı bir isim vermek için. \\
QML\_SINGLETON & Tek örnekli (singleton) tip. Bölüm 8'e bakın. \\
\rowcolor{acikGri} QML\_UNCREATABLE("msg") & QML'de \kod{new}lenemez ama tipi/enum'u görünür. \\
QML\_ANONYMOUS & İsimsiz; sadece özellik tipi olarak döner, QML'de yaratılamaz. \\
\rowcolor{acikGri} QML\_FOREIGN & Değiştiremediğiniz (3. parti/Qt) bir tipi kaydetmek için. \\
\bottomrule
\end{tabularx}
\end{center}

\subsection{Eski Yöntem: qmlRegisterType (Qt 5 tarzı)}

Eski projelerde veya modül kullanmadığınızda kaydı \kod{main.cpp} içinden
elle yaparsınız. Bilmekte fayda var:

\begin{lstlisting}
#include "backend.h"
// ...
int main(int argc, char *argv[])
{
    QGuiApplication app(argc, argv);

    // qmlRegisterType<Sinif>("Uri", surumBuyuk, surumKucuk, "QmlAdi");
    qmlRegisterType<Backend>("MyApp", 1, 0, "Backend");

    QQmlApplicationEngine engine;
    engine.load(QUrl(QStringLiteral("qrc:/Main.qml")));
    return app.exec();
}
\end{lstlisting}

\begin{uyarikutu}
Yeni projelerde \kod{qmlRegisterType} yerine \kod{QML\_ELEMENT} + modül
yaklaşımını tercih edin. Elle kayıt; sürüm numaralarını manuel yönetmenizi
gerektirir, derleme zamanı kontrolü sağlamaz ve QML derleyici
optimizasyonlarından yararlanamaz.
\end{uyarikutu}

% =====================================================================
\section{Q\_PROPERTY: Veri Açmanın Kalbi}
% =====================================================================

\kod{Q\_PROPERTY}, bir C++ üye değişkenini QML'e \emph{özellik} olarak açar.
QML bunu okuyabilir, yazabilir ve — en önemlisi — değeri değiştiğinde arayüz
\textbf{otomatik olarak} güncellenir. Bu, QML'in \emph{property binding}
(özellik bağlama) gücünün C++ tarafındaki karşılığıdır.

\subsection{Anatomi: READ, WRITE, NOTIFY}

\begin{lstlisting}
Q_PROPERTY(tip ad READ okuyucu WRITE yazici NOTIFY degistiSinyali)
\end{lstlisting}

\begin{center}
\begin{tabularx}{\textwidth}{@{}>{\ttfamily\bfseries\color{vurguMavi}}l >{\raggedright\arraybackslash}X@{}}
\toprule
\rowcolor{acikMavi}\color{anaMavi}\textbf{Anahtar} & \color{anaMavi}\textbf{Açıklama} \\
\midrule
READ & Değeri döndüren getter metodu. Zorunlu (MEMBER yoksa). \\
\rowcolor{acikGri} WRITE & Değeri ayarlayan setter. Opsiyonel — yoksa salt-okunur. \\
NOTIFY & Değer değişince yayılan sinyal. Binding güncellemesi için \textbf{şart}. \\
\rowcolor{acikGri} CONSTANT & Değer hiç değişmez; NOTIFY gerekmez. \\
MEMBER & Getter/setter yazmadan üyeyi doğrudan açar. \\
\rowcolor{acikGri} REQUIRED & QML nesnesi oluşturulurken bu özellik atanmak zorunda. \\
\bottomrule
\end{tabularx}
\end{center}

\subsection{Tam Örnek: Sayaç Sınıfı}

\begin{ornekkutu}[counter.h]
\end{ornekkutu}
\begin{lstlisting}
#ifndef COUNTER_H
#define COUNTER_H

#include <QObject>
#include <QQmlEngine>

class Counter : public QObject
{
    Q_OBJECT
    QML_ELEMENT

    // "count" adinda, okunabilir/yazilabilir bir int ozellik
    Q_PROPERTY(int count READ count WRITE setCount NOTIFY countChanged)

public:
    explicit Counter(QObject *parent = nullptr) : QObject(parent) {}

    int count() const { return m_count; }

    void setCount(int value)
    {
        if (m_count == value)   // Deger degismediyse cik (sonsuz donguyu onler)
            return;
        m_count = value;
        emit countChanged();    // QML'e "degistim" de
    }

    Q_INVOKABLE void increment() { setCount(m_count + 1); }
    Q_INVOKABLE void reset()     { setCount(0); }

signals:
    void countChanged();        // NOTIFY sinyali

private:
    int m_count = 0;
};

#endif
\end{lstlisting}

\begin{uyarikutu}[En kritik iki kural]
\textbf{1.} Setter içinde \emph{daima} değer değişti mi diye kontrol edin. Bu
kontrol olmadan binding'ler kendini tetikleyip \textbf{sonsuz döngüye} girebilir.\\
\textbf{2.} Değeri değiştiren her yolda NOTIFY sinyalini \kod{emit} edin. Etmezseniz
QML arayüzü \emph{sessizce eski değeri gösterir} ve saatlerce hata ararsınız.
\end{uyarikutu}

\subsection{QML Tarafı: Binding Sihiri}

\begin{lstlisting}[style=qmlstyle]
import QtQuick
import QtQuick.Controls
import MyApp

Window {
    width: 300; height: 200; visible: true

    Counter { id: counter }

    Column {
        anchors.centerIn: parent
        spacing: 12

        // counter.count degistiginde bu Text OTOMATIK guncellenir
        Text {
            text: "Sayac: " + counter.count
            font.pixelSize: 24
        }

        Button {
            text: "Arttir (+1)"
            onClicked: counter.increment()   // C++ metodu cagrildi
        }
        Button {
            text: "Sifirla"
            onClicked: counter.count = 0      // WRITE ile dogrudan atama
        }
    }
}
\end{lstlisting}

\begin{ipucukutu}
Burada \kod{text: "Sayac: " + counter.count} bir \textbf{binding}'dir. QML motoru
\kod{count}'un \kod{countChanged} sinyaline otomatik abone olur; sinyal her
yayıldığında \kod{Text} kendini yeniler. Siz elle güncelleme kodu yazmazsınız —
bildirimsel (declarative) programlamanın özü budur.
\end{ipucukutu}

\subsection{Kısayol: MEMBER ile Hızlı Özellik}

Getter/setter yazmak istemediğiniz basit durumlarda \kod{MEMBER} kullanın:

\begin{lstlisting}
class Settings : public QObject
{
    Q_OBJECT
    QML_ELEMENT
    // Getter/setter yok; uyeye dogrudan baglanir. NOTIFY yine onerilir.
    Q_PROPERTY(QString username MEMBER m_username NOTIFY usernameChanged)
public:
    explicit Settings(QObject *parent = nullptr) : QObject(parent) {}
signals:
    void usernameChanged();
private:
    QString m_username;
};
\end{lstlisting}

\subsection{Salt-Okunur ve Sabit Özellikler}

\begin{lstlisting}
class AppInfo : public QObject
{
    Q_OBJECT
    QML_ELEMENT
    // WRITE yok -> QML sadece okuyabilir
    Q_PROPERTY(QString status READ status NOTIFY statusChanged)
    // CONSTANT -> hic degismez, NOTIFY gerekmez
    Q_PROPERTY(QString version READ version CONSTANT)
public:
    explicit AppInfo(QObject *parent = nullptr) : QObject(parent) {}
    QString status() const  { return m_status; }
    QString version() const { return QStringLiteral("1.0.0"); }
signals:
    void statusChanged();
private:
    QString m_status = QStringLiteral("Hazir");
};
\end{lstlisting}

% =====================================================================
\section{Q\_INVOKABLE ve Slotlar: C++ Metotlarını Çağırmak}
% =====================================================================

QML'in bir C++ metodunu çağırabilmesi için metodun meta-object sistemine
kayıtlı olması gerekir. Bunun iki yolu vardır: \kod{Q\_INVOKABLE} ile
işaretlemek veya metodu bir \kod{slot} yapmak.

\subsection{İki Yöntem, Aynı Sonuç}

\begin{lstlisting}
class Calculator : public QObject
{
    Q_OBJECT
    QML_ELEMENT
public:
    explicit Calculator(QObject *parent = nullptr) : QObject(parent) {}

    // Yontem 1: Q_INVOKABLE ile isaretle
    Q_INVOKABLE double topla(double a, double b) { return a + b; }
    Q_INVOKABLE double bol(double a, double b)
    {
        if (b == 0.0)
            return 0.0;      // Sifira bolme korumasi
        return a / b;
    }

public slots:
    // Yontem 2: public slot -> QML'den otomatik cagrilabilir
    QString formatla(double sayi)
    {
        return QString::number(sayi, 'f', 2);   // "3.14" gibi
    }
};
\end{lstlisting}

\begin{lstlisting}[style=qmlstyle]
Calculator { id: calc }

Text {
    // Her ucu de QML'den dogrudan cagrilabilir
    text: calc.formatla(calc.bol(calc.topla(10, 5), 3))  // "5.00"
}
\end{lstlisting}

\begin{notkutu}[Q\_INVOKABLE mı, slot mu?]
İkisi de QML'den çağrılabilir. \textbf{Slot}; ayrıca bir sinyale
\kod{connect} edilebildiği için sinyal--slot mimarisine uygundur. \textbf{Q\_INVOKABLE}
ise "bu sadece QML'den çağrılan bir metot" niyetini daha net belirtir. Modern
tercih: sinyale bağlanmayacaksa \kod{Q\_INVOKABLE} kullanın.
\end{notkutu}

\subsection{Parametre ve Dönüş Tipleri}

QML ile C++ arasında otomatik dönüştürülen tipler (\kls{QVariant} üzerinden):

\begin{center}
\begin{tabularx}{\textwidth}{@{}>{\ttfamily\color{anaMavi}}l >{\ttfamily\color{mor}}l >{\raggedright\arraybackslash}X@{}}
\toprule
\rowcolor{acikMavi}\color{anaMavi}\textbf{C++ Tipi} & \color{anaMavi}\textbf{QML Karşılığı} & \color{anaMavi}\textbf{Not} \\
\midrule
int, double, bool & int, real, bool & Doğrudan \\
\rowcolor{acikGri} QString & string & UTF-16 \\
QVariant & var & Her şeyi taşır \\
\rowcolor{acikGri} QVariantList & list/array & JS dizisi \\
QVariantMap & object & JS nesnesi (anahtar-değer) \\
\rowcolor{acikGri} QObject* & (nesne) & Referans olarak geçer \\
QDateTime & date & Tarih/saat \\
\bottomrule
\end{tabularx}
\end{center}

\subsection{JavaScript Nesnesi Döndürmek}

QML'e yapılandırılmış veri döndürmenin pratik yolu \kls{QVariantMap}'tir:

\begin{lstlisting}
Q_INVOKABLE QVariantMap kullaniciBilgisi(int id)
{
    QVariantMap m;
    m["id"]    = id;
    m["ad"]    = QStringLiteral("Ayse");
    m["yas"]   = 30;
    m["aktif"] = true;
    return m;   // QML'de bir JS nesnesi olur
}
\end{lstlisting}

\begin{lstlisting}[style=qmlstyle]
var u = backend.kullaniciBilgisi(42)
console.log(u.ad, u.yas)     // "Ayse 30"
\end{lstlisting}

% =====================================================================
\section{Sinyaller: C++'tan QML'i Haberdar Etmek}
% =====================================================================

Özellikler "değer" iletmek için idealdir; ama "bir olay oldu" (indirme bitti,
hata oluştu, bağlantı koptu) demek için \textbf{sinyaller} kullanılır. C++'ta
sinyal yayarsınız, QML tarafında ona \emph{handler} bağlarsınız.

\subsection{C++ Tarafı: Sinyal Tanımlama ve Yayma}

\begin{lstlisting}
class Downloader : public QObject
{
    Q_OBJECT
    QML_ELEMENT
public:
    explicit Downloader(QObject *parent = nullptr) : QObject(parent) {}

    Q_INVOKABLE void baslat(const QString &url)
    {
        emit basladi(url);
        // ... indirme mantigi (basitlestirilmis) ...
        for (int p = 0; p <= 100; p += 25)
            emit ilerleme(p);         // Parametreli sinyal
        emit tamamlandi(url, true);   // Cok parametreli sinyal
    }

signals:
    void basladi(const QString &url);
    void ilerleme(int yuzde);
    void tamamlandi(const QString &url, bool basarili);
};
\end{lstlisting}

\subsection{QML Tarafı: on\textless Sinyal\textgreater\ Handler'ı}

Qt otomatik olarak her \kod{sinyalAdi} için \kod{onSinyalAdi} handler'ı üretir
(ilk harf büyür):

\begin{lstlisting}[style=qmlstyle]
Downloader {
    id: downloader

    // "basladi" -> "onBasladi"
    onBasladi: (url) => console.log("Basladi:", url)

    // "ilerleme" -> "onIlerleme"; parametre adi sinyaldeki ile ayni
    onIlerleme: (yuzde) => progressBar.value = yuzde / 100

    onTamamlandi: (url, basarili) => {
        if (basarili)
            statusText.text = "Bitti: " + url
        else
            statusText.text = "Hata!"
    }
}
\end{lstlisting}

\subsection{Connections: Dışarıdan Sinyale Bağlanma}

Nesnenin tanımıyla aynı yerde handler yazamıyorsanız (ör. context property ile
gelen bir nesne) \kls{Connections} kullanın:

\begin{lstlisting}[style=qmlstyle]
Connections {
    target: downloader          // Hangi nesnenin sinyalleri
    function onTamamlandi(url, basarili) {
        console.log("Connections ile yakalandi:", url)
    }
}
\end{lstlisting}

\begin{ipucukutu}
Sinyal parametrelerine handler içinde \textbf{isimleriyle} erişebilirsiniz. Eski
Qt 5 sözdiziminde bunlar sihirli biçimde kapsamdaydı; Qt 6'da modern ve önerilen
biçim, yukarıdaki gibi \textbf{ok fonksiyonu} veya \kod{function onX(param)}
parametre listesidir. Bu, isim çakışmalarını önler.
\end{ipucukutu}

% =====================================================================
\section{Enum'ları QML'e Açmak}
% =====================================================================

C++ enum'larını QML'de kullanmak, "sihirli sayılar" yerine anlamlı isimler
kullanmayı sağlar. \kod{Q\_ENUM} makrosu enum'u meta-object sistemine kaydeder.

\subsection{Q\_ENUM ile Kayıt}

\begin{lstlisting}
class Player : public QObject
{
    Q_OBJECT
    QML_ELEMENT
public:
    explicit Player(QObject *parent = nullptr) : QObject(parent) {}

    enum State {
        Stopped,
        Playing,
        Paused
    };
    Q_ENUM(State)     // <-- Enum'u QML'e acar

    Q_PROPERTY(State state READ state WRITE setState NOTIFY stateChanged)

    State state() const { return m_state; }
    void setState(State s) {
        if (m_state == s) return;
        m_state = s;
        emit stateChanged();
    }
signals:
    void stateChanged();
private:
    State m_state = Stopped;
};
\end{lstlisting}

\subsection{QML'de Enum Kullanımı}

Enum değerlerine \kod{TipAdi.DegerAdi} biçiminde erişilir:

\begin{lstlisting}[style=qmlstyle]
Player {
    id: player

    onStateChanged: {
        switch (state) {
        case Player.Playing: console.log("Oynatiliyor"); break
        case Player.Paused:  console.log("Duraklatildi"); break
        case Player.Stopped: console.log("Durduruldu");  break
        }
    }
}

Button {
    text: "Oynat"
    // Sayi yerine anlamli isim:
    onClicked: player.state = Player.Playing
}
\end{lstlisting}

\begin{uyarikutu}
Enum değerlerine QML'de \textbf{sınıf adı} üzerinden erişilir (\kod{Player.Playing}),
nesne adı üzerinden değil (\kod{player.Playing} \textbf{yanlıştır}). Ayrıca enum
kaydı için sınıfın QML'e kayıtlı olması (\kod{QML\_ELEMENT}) gerekir; sadece
enum'u açmak istiyor ama nesne oluşturulmasını istemiyorsanız
\kod{QML\_UNCREATABLE} ile birleştirin.
\end{uyarikutu}

% =====================================================================
\section{Context Property: Global Nesne Enjeksiyonu}
% =====================================================================

\kod{QML\_ELEMENT} ile tipleri açtınız — ama bazen QML'in \emph{belirli bir C++
örneğine} (tip değil, hazır nesne) erişmesini istersiniz. Örneğin veritabanına
bağlı tek bir \kod{backend} nesnesi. Bunun klasik yolu \textbf{context property}'dir.

\subsection{setContextProperty ile Nesne Enjekte Etme}

\begin{lstlisting}
#include <QQmlContext>
#include "backend.h"

int main(int argc, char *argv[])
{
    QGuiApplication app(argc, argv);
    QQmlApplicationEngine engine;

    Backend backend;   // C++'ta olusturulan tek ornek
    backend.setUser("admin");

    // "backend" adiyla TUM QML dosyalarina global olarak enjekte et
    engine.rootContext()->setContextProperty("backend", &backend);

    engine.loadFromModule("MyApp", "Main");
    return app.exec();
}
\end{lstlisting}

\begin{lstlisting}[style=qmlstyle]
// Herhangi bir QML dosyasinda, import bile gerekmeden:
Text {
    text: "Kullanici: " + backend.user   // dogrudan "backend"
}
\end{lstlisting}

\begin{uyarikutu}[Context property'nin dezavantajları]
Context property'ler pratik görünse de modern Qt'de \textbf{önerilmez}:
\textbf{(1)} Derleme zamanı tip kontrolü yoktur — yanlış özellik adı sadece
çalışma zamanında sessizce \kod{undefined} döner. \textbf{(2)} QML derleyici
(qmltc) bunları optimize edemez. \textbf{(3)} Her QML dosyasında görünür olması
kapsam kirliliği yaratır. Mümkünse yerine \textbf{singleton} (Bölüm 8) kullanın.
\end{uyarikutu}

\subsection{Alternatif: initialProperties ile Örnek Aktarımı}

Tek bir kök nesneye örnek geçirmek için daha temiz bir yol:

\begin{lstlisting}
Backend backend;
QQmlApplicationEngine engine;
engine.loadFromModule("MyApp", "Main");
// Kök nesnenin "backend" ozelligine C++ ornegini bagla
// (Main.qml icinde: property Backend backend)
\end{lstlisting}

% =====================================================================
\section{Singleton'lar: Uygulama Geneli Tek Nesne}
% =====================================================================

Uygulama boyunca tek bir örneği olması gereken servisler (ayarlar yöneticisi,
oturum, tema, veritabanı erişimi) için \textbf{singleton} idealdir. Context
property'nin tip-güvenli, optimize edilebilir modern alternatifidir.

\subsection{QML\_SINGLETON ile Tanımlama}

\begin{lstlisting}
class AppSettings : public QObject
{
    Q_OBJECT
    QML_ELEMENT
    QML_SINGLETON       // <-- Tek ornek olacak
    Q_PROPERTY(QString theme READ theme WRITE setTheme NOTIFY themeChanged)
    Q_PROPERTY(int fontSize READ fontSize WRITE setFontSize NOTIFY fontSizeChanged)
public:
    explicit AppSettings(QObject *parent = nullptr) : QObject(parent) {}

    QString theme() const { return m_theme; }
    void setTheme(const QString &t) {
        if (m_theme == t) return;
        m_theme = t; emit themeChanged();
    }
    int fontSize() const { return m_fontSize; }
    void setFontSize(int s) {
        if (m_fontSize == s) return;
        m_fontSize = s; emit fontSizeChanged();
    }
signals:
    void themeChanged();
    void fontSizeChanged();
private:
    QString m_theme = QStringLiteral("dark");
    int m_fontSize = 14;
};
\end{lstlisting}

\subsection{QML'de Singleton Kullanımı}

Singleton'a \textbf{tip adıyla} doğrudan erişilir; nesne oluşturmaya gerek yoktur:

\begin{lstlisting}[style=qmlstyle]
import QtQuick
import MyApp

Window {
    // Tum uygulamada ayni ornek. Nesne olusturmadan tip adiyla eris:
    color: AppSettings.theme === "dark" ? "#222" : "#fff"

    Text {
        text: "Ayarlar"
        font.pixelSize: AppSettings.fontSize    // Global, canli baglanti
        color: AppSettings.theme === "dark" ? "white" : "black"
    }

    Button {
        text: "Temayi Degistir"
        onClicked: AppSettings.theme =
                   AppSettings.theme === "dark" ? "light" : "dark"
    }
}
\end{lstlisting}

\begin{ipucukutu}
Bir singleton'ın \kod{theme} özelliğini bir yerde değiştirdiğinizde,
uygulamadaki \emph{tüm} binding'ler anında güncellenir. Bu, tema/dil/ayar gibi
global durumları yönetmenin en zarif yoludur. C++ tarafında singleton örneğine
\kod{engine} üzerinden de erişebilirsiniz.
\end{ipucukutu}

\subsection{Özel Kurucu Gereken Singleton}

Örneğin oluşumunu kontrol etmek isterseniz \kod{create} fonksiyonu kullanın:

\begin{lstlisting}
class Database : public QObject
{
    Q_OBJECT
    QML_ELEMENT
    QML_SINGLETON
public:
    // QML motoru bu statik fonksiyonu cagirarak ornegi olusturur
    static Database *create(QQmlEngine *, QJSEngine *engine)
    {
        Database *db = new Database();
        db->connectToDb();
        return db;   // Sahiplik QML motoruna gecer
    }
private:
    explicit Database(QObject *parent = nullptr) : QObject(parent) {}
    void connectToDb() { /* ... */ }
};
\end{lstlisting}

% =====================================================================
\section{Liste Modelleri: QAbstractListModel}
% =====================================================================

Tek nesne değil de bir \emph{liste} (kişiler, mesajlar, ürünler) göstermek
istediğinizde, C++ verisini \kls{ListView}'e bağlamanın performanslı yolu
\kls{QAbstractListModel} türetmektir.

\subsection{Rol Tabanlı Model Mantığı}

QML modelleri \textbf{rol (role)} kavramıyla çalışır. Her satır birden çok
alana sahiptir (ad, e-posta, yaş); her alan bir "rol"dür. QML delegate'i bu
rollere isimle erişir.

\begin{ornekkutu}[contactmodel.h]
\end{ornekkutu}
\begin{lstlisting}
#ifndef CONTACTMODEL_H
#define CONTACTMODEL_H

#include <QAbstractListModel>
#include <QQmlEngine>

struct Contact {
    QString name;
    QString email;
};

class ContactModel : public QAbstractListModel
{
    Q_OBJECT
    QML_ELEMENT
public:
    // Rolleri tanimla (Qt::UserRole'dan baslar)
    enum Roles {
        NameRole = Qt::UserRole + 1,
        EmailRole
    };

    explicit ContactModel(QObject *parent = nullptr) : QAbstractListModel(parent)
    {
        m_data = {
            {"Ali",  "ali@ornek.com"},
            {"Veli", "veli@ornek.com"}
        };
    }

    // 1) Kac satir var?
    int rowCount(const QModelIndex &parent = QModelIndex()) const override
    {
        Q_UNUSED(parent)
        return m_data.count();
    }

    // 2) Belirli satir+rol icin veri
    QVariant data(const QModelIndex &index, int role) const override
    {
        if (!index.isValid() || index.row() >= m_data.count())
            return {};
        const Contact &c = m_data.at(index.row());
        switch (role) {
        case NameRole:  return c.name;
        case EmailRole: return c.email;
        }
        return {};
    }

    // 3) Rol numaralarini QML'deki isimlere esle
    QHash<int, QByteArray> roleNames() const override
    {
        return {
            { NameRole,  "name"  },
            { EmailRole, "email" }
        };
    }

    // Satir eklemek: beginInsertRows/endInsertRows sart!
    Q_INVOKABLE void addContact(const QString &name, const QString &email)
    {
        beginInsertRows(QModelIndex(), m_data.count(), m_data.count());
        m_data.append({ name, email });
        endInsertRows();
    }

private:
    QList<Contact> m_data;
};

#endif
\end{lstlisting}

\begin{uyarikutu}[Üç zorunlu metot + değişiklik bildirimi]
\kod{rowCount}, \kod{data} ve \kod{roleNames} \textbf{override edilmek zorundadır}.
Ayrıca veriyi değiştirirken \kod{beginInsertRows}/\kod{endInsertRows} (ekleme),
\kod{beginRemoveRows}/\kod{endRemoveRows} (silme) veya \kod{dataChanged} (güncelleme)
çağrılarını \textbf{atlamayın}. Atlarsanız \kls{ListView} güncellenmez veya
çökme yaşarsınız.
\end{uyarikutu}

\subsection{QML'de ListView ile Gösterme}

\begin{lstlisting}[style=qmlstyle]
import QtQuick
import QtQuick.Controls
import MyApp

ListView {
    anchors.fill: parent
    model: ContactModel {}    // C++ modeli dogrudan bagla

    delegate: Rectangle {
        width: ListView.view.width
        height: 50
        color: index % 2 ? "#f0f0f0" : "white"

        Row {
            spacing: 10
            anchors.verticalCenter: parent.verticalCenter
            // roleNames'teki isimler dogrudan kapsamda:
            Text { text: name; font.bold: true }
            Text { text: email; color: "gray" }
        }
    }
}
\end{lstlisting}

\begin{notkutu}[Basit listeler için kısayol]
Sadece birkaç öğelik statik/basit bir liste için \kls{QAbstractListModel}
fazla ağır kalır. \kod{QStringList} veya \kod{QVariantList} döndüren bir
\kod{Q\_PROPERTY} de \kod{model} olarak kullanılabilir. Ancak büyük, sık
değişen veya çok alanlı listeler için mutlaka \kls{QAbstractListModel} tercih
edin — performans farkı büyüktür.
\end{notkutu}

% =====================================================================
\section{C++'tan QML'e Erişmek}
% =====================================================================

Şimdiye kadar akış çoğunlukla QML $\rightarrow$ C++ yönündeydi. Bazen tersine,
C++'tan bir QML nesnesinin özelliğini okumak, metodunu çağırmak veya sinyaline
bağlanmak istersiniz. Bu genelde \emph{gerekli olmadıkça kaçınılması} gereken
bir yöndür ama gerektiğinde şöyle yapılır.

\subsection{Kök Nesneye ve Alt Nesnelere Erişim}

\begin{lstlisting}
QQmlApplicationEngine engine;
engine.loadFromModule("MyApp", "Main");

// Kok QML nesnesi (Window)
QObject *root = engine.rootObjects().first();

// Ozellik oku/yaz (property adiyla)
root->setProperty("title", "C++ tarafindan degistirildi");
QVariant w = root->property("width");

// objectName ile bir alt nesne bul
//   (QML'de: Text { objectName: "durumMetni" })
QObject *label = root->findChild<QObject*>("durumMetni");
if (label)
    label->setProperty("text", "Guncellendi");
\end{lstlisting}

\subsection{C++'tan QML Fonksiyonu Çağırmak}

QML'de \kod{function} olarak tanımlı bir metodu \kls{QMetaObject::invokeMethod}
ile çağırabilirsiniz:

\begin{lstlisting}[style=qmlstyle]
// Main.qml
Window {
    function mesajGoster(metin) {
        console.log("QML fonksiyonu:", metin)
        return "alindi: " + metin
    }
}
\end{lstlisting}

\begin{lstlisting}
QObject *root = engine.rootObjects().first();
QVariant donus;
QMetaObject::invokeMethod(root, "mesajGoster",
    Q_RETURN_ARG(QVariant, donus),
    Q_ARG(QVariant, QString("Selam QML")));
qDebug() << donus;   // "alindi: Selam QML"
\end{lstlisting}

\begin{uyarikutu}
C++'tan QML nesnelerine \kod{findChild}/\kod{invokeMethod} ile erişmek
\textbf{kırılgan} bir desendir: QML yapısını değiştirdiğinizde C++ kodu sessizce
bozulur ve derleyici sizi uyarmaz. \textbf{Tercih edilen mimari}: veriyi ve
mantığı C++'ta tutun, \kod{Q\_PROPERTY} ve sinyallerle QML'e \emph{itin}; QML
bunlara binding ile tepki versin. C++'ın QML'e uzanması yerine QML'in C++'a
bağlanması daha sağlıklıdır.
\end{uyarikutu}

% =====================================================================
\section{Değer Tipleri ve Gadget'lar (Q\_GADGET)}
% =====================================================================

Her veri yapısının tam bir \kls{QObject} olması gerekmez. Konum, renk, boyut
gibi küçük, kopyalanabilir veri taşıyıcıları için \kod{Q\_GADGET} kullanılır —
sinyal/slot yükü olmadan meta bilgi sağlar.

\begin{lstlisting}
class Point3D
{
    Q_GADGET     // QObject degil! Hafif meta-tip
    Q_PROPERTY(double x READ x WRITE setX)
    Q_PROPERTY(double y READ y WRITE setY)
    Q_PROPERTY(double z READ z WRITE setZ)
public:
    double x() const { return m_x; }  void setX(double v) { m_x = v; }
    double y() const { return m_y; }  void setY(double v) { m_y = v; }
    double z() const { return m_z; }  void setZ(double v) { m_z = v; }

    Q_INVOKABLE double length() const
    {
        return std::sqrt(m_x*m_x + m_y*m_y + m_z*m_z);
    }
private:
    double m_x = 0, m_y = 0, m_z = 0;
};
\end{lstlisting}

\begin{center}
\begin{tabularx}{\textwidth}{@{}>{\bfseries\color{anaMavi}}l >{\raggedright\arraybackslash}X >{\raggedright\arraybackslash}X@{}}
\toprule
\rowcolor{acikMavi}\color{anaMavi}\textbf{Özellik} & \color{anaMavi}\textbf{Q\_OBJECT} & \color{anaMavi}\textbf{Q\_GADGET} \\
\midrule
Sinyal/slot & Var & Yok \\
\rowcolor{acikGri} Kopyalanabilir & Hayır (kimliği var) & Evet (değer tipi) \\
Bellek yükü & Yüksek & Düşük \\
\rowcolor{acikGri} Kullanım & Servis, model, mantık & Küçük veri taşıyıcı \\
\bottomrule
\end{tabularx}
\end{center}

% =====================================================================
\section{Yaygın Hatalar ve En İyi Pratikler}
% =====================================================================

\subsection{Sık Yapılan Hatalar}

\begin{center}
\begin{tabularx}{\textwidth}{@{}>{\raggedright\arraybackslash\bfseries\color{kirmizi}}p{5cm} >{\raggedright\arraybackslash}X@{}}
\toprule
\rowcolor{kirmizi!12}\color{kirmizi}\textbf{Hata} & \color{kirmizi}\textbf{Sonuç \& Çözüm} \\
\midrule
NOTIFY sinyalini emit etmemek & QML eski değeri gösterir. Setter'da değer değişince mutlaka \kod{emit} edin. \\
\rowcolor{acikGri} Setter'da değişiklik kontrolü yapmamak & Binding sonsuz döngüye girer. \kod{if (m == v) return;} ekleyin. \\
Q\_OBJECT makrosunu unutmak & Sinyaller/özellikler çalışmaz, linker hatası. Sınıf başına ekleyin. \\
\rowcolor{acikGri} Argümansız kurucu olmaması & QML nesneyi oluşturamaz. Ya kurucu ekleyin ya \kod{QML\_UNCREATABLE}. \\
Modeli \kod{begin/end...} olmadan değiştirmek & ListView bozulur/çöker. Değişiklik metotlarını sarmalayın. \\
\rowcolor{acikGri} Nesne sahipliğini yanlış yönetmek & Çift silme veya sızıntı. QML'e verdiğiniz nesnelerde JS sahipliğine dikkat. \\
\bottomrule
\end{tabularx}
\end{center}

\subsection{Nesne Sahipliği (Ownership)}

\begin{bilgikutu}
Bir C++ metodu QML'e \kod{QObject*} döndürdüğünde sahiplik kuralı: nesnenin
\textbf{parent'ı varsa} C++ sahiptir; \textbf{parent'ı yoksa} QML çöp toplayıcısı
(garbage collector) onu silebilir. QML'in beklenmedik anda silmesini istemiyorsanız
nesneye bir parent verin veya \kod{QQmlEngine::setObjectOwnership(obj,
QQmlEngine::CppOwnership)} ile sahipliği açıkça C++'a bırakın.
\end{bilgikutu}

\subsection{Mimari Tavsiyeleri}

\begin{itemize}[leftmargin=1.4em]
\item \textbf{Tek yön akış:} Veri ve mantık C++'ta; QML sadece \emph{sunum}.
      C++ QML'e veri iter, QML C++ metotlarını çağırır. C++'ın QML iç yapısına
      uzanmasından kaçının.
\item \textbf{Singleton \textgreater\ context property:} Global servisler için
      tip-güvenli, optimize edilebilir singleton'ları tercih edin.
\item \textbf{Modül kullanın:} \kod{qt\_add\_qml\_module} + \kod{QML\_ELEMENT},
      elle \kod{qmlRegisterType}'a yeğdir.
\item \textbf{Ağır işi bloklamayın:} Uzun süren C++ işlemlerini ayrı bir iş
      parçacığında (\kls{QThread}, \kls{QtConcurrent}) yapın; bitince sinyalle
      QML'i haberdar edin. Aksi hâlde arayüz donar.
\item \textbf{const \& referans:} Metot parametrelerinde büyük tipleri
      \kod{const QString\&} gibi referansla alın; gereksiz kopyalamayı önler.
\end{itemize}

\subsection{Hızlı Referans Tablosu}

\begin{center}
\begin{tabularx}{\textwidth}{@{}>{\ttfamily\color{anaMavi}\small}l >{\raggedright\arraybackslash}X@{}}
\toprule
\rowcolor{acikMavi}\color{anaMavi}\textbf{Makro / Öğe} & \color{anaMavi}\textbf{Amaç} \\
\midrule
Q\_OBJECT & Meta-object desteği; her açılan sınıfta zorunlu. \\
\rowcolor{acikGri} QML\_ELEMENT & Sınıfı QML'e kaydeder. \\
QML\_SINGLETON & Tek örnekli tip. \\
\rowcolor{acikGri} QML\_UNCREATABLE & Görünür ama QML'de yaratılamaz (enum açmak için ideal). \\
Q\_PROPERTY & Veriyi özellik olarak açar (READ/WRITE/NOTIFY). \\
\rowcolor{acikGri} Q\_INVOKABLE & Metodu QML'den çağrılabilir yapar. \\
Q\_ENUM & Enum'u QML'e açar. \\
\rowcolor{acikGri} Q\_GADGET & Hafif, kopyalanabilir değer tipi. \\
signals: & Sinyal bölümü; C++'tan QML'e olay bildirimi. \\
\bottomrule
\end{tabularx}
\end{center}

\vspace{0.6cm}
\begin{tcolorbox}[colback=qtYesil!8,colframe=qtYesil!70!koyuGri,arc=3pt,boxrule=1pt]
\centering
\textbf{Özet:} QML arayüzü, C++ mantığı. İkisini \kod{Q\_PROPERTY} (veri),
\kod{Q\_INVOKABLE} (metot), \textbf{signal} (olay) üçlüsüyle birbirine bağlayın.
Tipleri \kod{QML\_ELEMENT} ile modül içinde kaydedin, global servisler için
singleton kullanın, listeler için \kls{QAbstractListModel} türetin. Akışı tek
yönlü tutun: \emph{C++ iter, QML bağlanır.}
\end{tcolorbox}

\end{document}
